\section{Fixed progressions and the Liouville weight}\label{sec:parity}

\subsection{Two residue conditions before division by shared labels}

Fix $M$ and a residue $r\pmod M$, which may be a nonunit. Multiply
the first physical vector by $\ind_{n\equiv r\ (M)}$ and the second
by $\ind_{n'\equiv r+1\ (M)}$. Their norms decrease. The first vector
remains independent of its lists at a fixed position. Thus the general
pairing still bounds the entire residual family.

For a full small prime interval $\mathcal P_g$, the harmonic law
$1/(V_gp)$ is uniform on $(\Z/M\Z)^\times$ to any fixed negative
power of $L$. Partial summation of the fixed-modulus prime estimate
proves this, since $\log p\geq L^{.27}$ and $V_g$ is bounded below.
One shared label from this group suffices. Conditional on every other
label in the independent-label model, multiplication by their product
permutes the units. The collision errors described in
Section~\ref{sec:framework} remain negligible.

Perform this averaging only after the dyads and Mellin integral have
been restored in \eqref{eq:strip}. For a unit $d\pmod M$, the exact
identity is
\begin{equation}\label{eq:orbit}
 \frac1{\varphi(M)}\sum_{t\in(\Z/M\Z)^\times}
 \ind_{tw\equiv r\ (M)}\ind_{t(w+d)\equiv r+1\ (M)}
 =\frac1{\varphi(M)}\ind_{w\equiv rd\ (M)}.
\end{equation}
Subtracting the two congruences forces $td\equiv1$. The unique
possible $t=d^{-1}$ satisfies the first congruence exactly when
$w\equiv rd$. There is no division by $r$ or $w$.
The raw pattern has $d=D_C=1$. Each comparison acquires
$w\equiv rD_C\pmod M$.

To bound those comparisons, first split each free product into its
fixed unit residues modulo $M$. The remaining condition on $w$ is a
fixed residue. With the Fourier convention
$\widehat U(\theta)=\sum_n U(n)\e(\theta n)$,
\[
 \widehat{\ind_{n\equiv s\ (M)}U}(\theta)
  =\frac1M\sum_{j\bmod M}\e(-js/M)\widehat U(\theta+j/M).
\]
Its squared energy is bounded by the corresponding average of the
squared energies. A rational center with denominator $k$ moves to
one with denominator at most $Mk$. Increasing the fixed logarithmic
denominator and width exponents covers all these arcs. In
\eqref{eq:scales}, $H'$ grows faster than every fixed power of $L$
times a small-prime scale, so the new time cutoff is still admissible.
The bounded tests on the free products preserve the bilinear multiplier
estimate.

Apply this argument to the source centered endpoint
$\W(T_x-C_x)$. The source's local energy estimate already holds at
every fixed logarithmic denominator and time exponent. It yields
\begin{equation}\label{eq:unsigned-correlation}
 \sum_{w\equiv r\ (M)}\frac{\Psi(w/x)}w
 \W(w)(T_x(w)-C_x)G(w+1)W_1(w)W_1(w+1)\ll_A L^{-A}.
\end{equation}
We will use this to prove the unsigned progression law. Placing the
residue condition inside $G$ would violate $G(n)=G(n_*)$; the two
physical residue conditions avoid that problem.

\subsection{Ordinary-integer cancellation}

The classical Liouville estimate, in the reduced-residue form recorded
by Pintz \cite[equation (1.5)]{Pintz}, is
\begin{equation}\label{eq:bvlambda}
 \sum_{k\leq R^\theta}\max_{(a,k)=1}
 \left|\sum_{\substack{n\leq R\\n\equiv a\ (k)}}\lambda(n)\right|
 \ll_{A,\theta}R(\log R)^{-A},\qquad \theta<\tfrac12.
\end{equation}
All parameters are fixed. Smooth compactly supported dyadic weights
follow by partial summation, with the sum over moduli inside the
integral of absolute cumulative sums. No supremum over endpoints is
inserted into \eqref{eq:bvlambda}.

If a residue $a\pmod k$ is a nonunit, set $g=(a,k)$ and write
$n=gv$. Then $v\equiv a/g\pmod{k/g}$ is a reduced class and
$\lambda(gv)=\lambda(g)\lambda(v)$. For
$R\geq x^\delta$, $k\leq L^C$, and $|t|\leq L^B$, this gives
\begin{equation}\label{eq:swlambda}
 \frac1R\sum_{n\equiv a\ (k)}\lambda(n)\psi(n/R)n^{it}
 \ll_A L^{-A}.
\end{equation}
The test $\psi$ has fixed compact support in $(0,\infty)$, with
fixed logarithmic derivative costs. The reduced scale $R/g$ is still
a positive power of $x$; increasing the accuracy in
\eqref{eq:bvlambda} pays every displayed logarithmic cost.

\subsection{A signed correlation with a nonempty tuple}

\begin{lemma}\label{lem:lambda-correlation}
For the marked data of Section~\ref{sec:framework}, fixed $M,r$, and
$\Psi\in C_c^\infty((0,\infty))$,
\begin{equation}\label{eq:lambda-correlation}
 \sum_{w\equiv r\ (M)}\frac{\Psi(w/x)}w
 \lambda(w)\W(w)T_x(w)G(w+1)W_1(w)W_1(w+1)
 \ll_A L^{-A}.
\end{equation}
The tuple dimension satisfies $d\geq1$.
\end{lemma}

\begin{proof}
Insert $\lambda(n)$ in the first physical vector. Its norm and
list independence are unchanged. If $s_P$ is the number of active
groups and $c_\nu=\sum_g|C_g|$, then
\[
 \Omega(D_R)=Js_P-c_\nu,\qquad
 \lambda(D_Rw)=(-1)^{Js_P}(-1)^{c_\nu}\lambda(w).
\]
These identities count repeated auxiliary labels with multiplicity.
Divide out the common sign $(-1)^{Js_P}$. The transference result
and \eqref{eq:orbit} give the raw correlation plus comparisons
multiplied by $(-1)^{c_\nu}$, up to $O_A(L^{-A})$.
Each comparison still needs a separate bound.

At low Mellin times $|t|\leq L^B$, expand the tuple and the divisor
model \eqref{eq:divisor}, writing $w=\ell av$.
The scale of $v$ is $R=Y/(\ell a)\geq x^{\eps-o(1)}$.
Complete multiplicativity separates $\lambda(\ell a)$, and splitting
the additive phase into residue classes permits \eqref{eq:swlambda}.
The total normalized error is bounded using
$\sum_{\boldsymbol\ell}1/\ell=O(1)$ and
$\sum|A_a|/a\leq L^C$. Cauchy's inequality controls the divisor-model
error without changing its square norm. Choose the accuracy after
$B$ so that the integral over these low times also saves every
required power of $L$.

At high times remove one small mark and one long prime slot, giving
$w=p_s p_\ell v$. Then
$\lambda(w)=\lambda(v)$, so the small- and long-prime polynomials
are unchanged up to constant signs. Their residual coefficient
has the same square norm as the original coefficient. Small-mark
squares and repeated numerical long-slot primes have square errors
$O(YL^C\exp(-L^{.27}))$ and $O(YL^Cx^{-\eps})$. Ordinary Dirichlet
mean square bounds dispose of these errors on the whole time interval.

The factorial estimate in \cite[Section 6]{D} places the times when
the small-prime polynomial exceeds a prescribed negative power of
$L$ inside at most $\exp(O(L^{.9}))$ unit intervals. Outside these
intervals, ordinary mean square on the collapsed scale $Y/P_s$
applies because $T_0/(Y/P_s)=LP_s/H\to0$. Here $P_s$ denotes the
small-prime scale, so $P_s\leq\exp(L^{.36})$.
On the exceptional intervals, use \eqref{eq:longprime} and the
source sparse Gram estimate for the remaining coefficients. That
estimate is coefficient-independent: its matrix entries are the
unweighted sums $\sum_{v\asymp R}v^{i(t_i-t_j)}$. Multiplying the
tested coefficient by $\lambda(v)$ preserves its square norm.
The long-prime scale lies between $x^\eps/2$ and $x^{1-\eps}$;
the residual scale is at least $x^{\eps-o(1)}$. All retained times
are below $x^2$ after fixed logarithmic twists. These are the source
estimate's ranges. Choose $B$ and the long-prime accuracy after the
fixed comparison costs.

This proves \eqref{eq:energy}. The Fourier shifts for the fixed
residue condition use the same bounds with enlarged fixed exponents.
The local Mellin and multiplier inputs now bound every comparison.
Subtract them from the transferred residual to obtain
\eqref{eq:lambda-correlation}.
\end{proof}

The nonempty tuple supplies the long prime in this proof. We have
not proved a signed constant-endpoint estimate by this method.
The total Liouville mean will follow instead from size-biased sampling.

\subsection{Progression remainders and scalar replacement}

Fix $0<c_0<\eps/8$ and a smooth nonnegative compact cutoff $\Phi$.
For $H$ equal to a fixed inert weight, a product of two such weights,
a constant, or a bounded linear combination of these, the source
represented-divisor expansion and CRT give
\begin{align}\label{eq:unsigned-progression}
 \sum_{\substack{r\leq x^{c_0}\\(r,M)=1}}
 \left|\sum_{\substack{r\mid u+1\\u+1\equiv b\ (M)}}
 \Phi(u/x)H(u)T_x(u)
 -\frac{x}{rM}\left(\int\Phi\right)C_x\E_{\rm div}H\right|
 \ll_A xL^{-A}.
\end{align}
The statement also holds with $T_x,C_x$ replaced by $1,1$.
Indeed every retained tuple/divisor product is eventually coprime
to fixed $M$. A divisor sharing a prime with $r$ gives no solutions;
otherwise CRT gives one class modulo $rM$. The source comparison of
the forbidden inert primes with its independent model has exponentially
small error. The total rounding error is
$x^{1-\eps+c_0+o(1)}$. Divisors with $(r,M)>1$ contribute zero.

The corresponding signed remainder is
\begin{equation}\label{eq:signed-progression}
 \sum_{r\leq x^{c_0}}
 \left|\sum_{\substack{r\mid u+1\\u+1\equiv b\ (M)}}
    \Phi(u/x)H(u)\lambda(u)T_x(u)\right|
 \ll_A xL^{-A}.
\end{equation}
To see this, write $u=\ell av$ in each retained divisor term.
When compatible, the congruences specify a class of $v$ modulo $rM$.
Its gcd with $rM$ is the fixed $g=(b-1,M)$. After $v=gt$, the
class is reduced modulo $rM/g$. The moduli are distinct as $r$ varies
and are less than $(x/(\ell ag))^{1/3}$ for large $x$.
Apply \eqref{eq:bvlambda} at that residual scale and sum with the
reciprocal tuple and divisor coefficients. For the omitted terms use
the positive first-omitted Bonferroni divisor envelope and CRT. Its
factorial reciprocal tail saves every logarithmic power, and its
rounding error has the bound above. This direct progression argument
avoids multiplying an $L^2$ error by $x^{c_0}$.

Finally, the scalar replacement extends to all unit-modulus weights
$f_z(n)=z^{\Omega(n)}$ used here. Selecting one prime from each of
$K$ inert bands and writing $mn-1=ah$ gives
\begin{equation}\label{eq:detphase}
 f_z(mn-1)F_0(mn-1)=z^K f_z(h)F_0(h)
\end{equation}
for an invariant bounded $F_0$. The factor on the right lies outside
the original Cauchy square and has modulus at most one. The scalar
coefficients $\alpha,\beta$ are unchanged, so the assumed coefficient
criterion applies. The selected-prime-square error is absolute.
The product class is imposed by
\[
 \ind_{mn\equiv b\ (M)}
 =\frac1{\varphi(M)}\sum_{\chi\bmod M}
       \overline{\chi(b)}\chi(m)\chi(n).
\]
The character factors preserve \eqref{eq:coefficient}, with logarithmic
test moduli enlarged by a fixed factor. No Liouville or complex weight
is inserted into the prime-minus-rough discrepancy coefficient.
